% Lab 10: Spacetime diagrams.
%
% One source, three PDFs:
%   Labs/Lab10_spacetime.tex        grayscale handout (this file): the moving
%                                   grid is dark and dashed, for B&W copies
%   Labs/Lab10_spacetime_color.tex  colour handout: the moving grid is blue
%   Labs/Lab10_spacetime_key.tex    answer key, in colour so the red answers
%                                   stand out from the grids
% The wrappers set a flag and \input this file, so edit only here.
%
% Layout. Every page break is chosen by hand with \newpage, and answer space
% is stretchable: \room{n} is a small floor plus n shares of whatever the page
% has left over. So each page fills to the bottom, a bigger n gets a bigger
% share, and a question is never left at the foot of a page without its space.
% If an edit makes a page overflow, move a \newpage rather than shrinking the
% numbers.
\documentclass{handout}
\usepackage{spacetime}
\ifdefined\labkey \printanswers \fi

\newcommand{\room}[1]{0.35in plus #1fill}

\assignment{Lab 10\ifprintanswers\ (Key)\fi}
\title{Spacetime Diagrams}
\duedate{}

\begin{document}
\maketitle

\section*{Reading the diagrams}

Throughout this lab, $c=1$. Distances are in light-years (ly) and times in
years (yr), so light covers 1~ly every 1~yr, and a light ray is a line at
$45^\circ$.

Every diagram prints two grids.
\begin{itemize}[nosep]
  \item The \textbf{gray square grid} is the Earth frame, with coordinates $(x,t)$.
  \item The \textbf{\ifdefined\spacetimecolor blue\else dark dashed\fi\ grid} is a
        rocket moving at $\beta = v/c = 0.6$, with coordinates $(x',t')$. Its
        $t'$ axis is the rocket's own worldline. Lines parallel to the $t'$ axis
        are places of constant $x'$, and lines parallel to the $x'$ axis are the
        rocket's lines of constant $t'$, its ``now.''
  \item Each step of the tilted grid is 1~ly or 1~yr \emph{in the rocket
        frame}, even though it looks longer on the page. To read an event's
        rocket coordinates, count tilted lines; do not measure with a ruler.
\end{itemize}
At this speed $\gamma = 1/\sqrt{1-\beta^2} = 1.25$.

%%%%%%%%%%%%%%%%%%%%%%%%%%%%%%%%%%%%%%%%%%%%%%%%%%%%%%%%%%%%%%%%%%%%%%%%%%%
\section{Building the diagram}
%%%%%%%%%%%%%%%%%%%%%%%%%%%%%%%%%%%%%%%%%%%%%%%%%%%%%%%%%%%%%%%%%%%%%%%%%%%

\begin{center}
\begin{spacetime}[0.8]{-2}{6}{0}{7}
  \movinggrid{0.6}{0}{0}
  \answer{
    \draw[st ans] (2,0) -- (2,7) node[pos=0.9, st anslabel, right] {lamp};
    \draw[st ans] (0,0) -- (4.2,7) node[pos=0.85, st anslabel, left] {rocket};
    \draw[st light] (0,0) -- (6,6);
    \draw[st light] (0,0) -- (-2,2);
    \node[st event, label={[st anslabel]right:$E$}] at (3,5) {};
    \node[st event, label={[st anslabel]right:$F$}] at (5,5) {};
  }
\end{spacetime}
\end{center}

\begin{questions}
  \question On the diagram, draw the worldline of each of these.
  \begin{parts}
    \part A lamp sitting still at $x = 2$~ly.
    \part A rocket that leaves $x=0$ at $t=0$ moving at $0.6c$. Where is it at $t=5$~yr?
    \part A flash of light sent both ways from $x = 0$ at $t = 0$.
  \end{parts}
  % answered on the diagram, so no floor: just what is left of page 1
  \begin{solution}[0pt plus 1fill]
    (a) A vertical line at $x=2$. (b) A straight line from the origin through
    $(x,t) = (3,5)$: it covers 0.6~ly each year. (c) Two $45^\circ$ lines from
    the origin, one each way.
  \end{solution}

\newpage
  \question Explain why no object's worldline can tilt farther from vertical
  than $45^\circ$.
  \begin{solution}[\room{1}]
    The tilt from vertical measures speed: $\Delta x/\Delta t$. A line beyond
    $45^\circ$ covers more than 1~ly per year, which is faster than light.
  \end{solution}

  \question Your rocket's worldline lies on top of one of the tilted lines.
  Which one? Explain why a rocket's own time axis has to be its own worldline.
  \begin{solution}[\room{1}]
    The $t'$ axis. In its own frame the rocket sits still at $x'=0$, and the
    set of all events at $x'=0$ is exactly the $t'$ axis.
  \end{solution}

  \question Mark the event $E$ at $(x,t) = (3,5)$ and the event $F$ at
  $(x,t) = (5,5)$. Count tilted lines to find each event's rocket coordinates
  $(x',t')$.
  \begin{solution}[\room{1}]
    $E$: $(x',t') = (0,4)$, on the rocket's own worldline.
    $F$: $(x',t') = (2.5, 2.5)$.
  \end{solution}

  \question $F$ lies on your light ray. What do the rocket's coordinates of
  $F$ say about the speed of light in the rocket frame?
  \begin{solution}[\room{1}]
    $x' = t'$, so in the rocket frame the light also covers 1~ly per year:
    both frames measure the same speed of light. That is the whole reason the
    tilted grid has to tilt symmetrically about the $45^\circ$ line.
  \end{solution}
\end{questions}

\newpage
%%%%%%%%%%%%%%%%%%%%%%%%%%%%%%%%%%%%%%%%%%%%%%%%%%%%%%%%%%%%%%%%%%%%%%%%%%%
\section{Relativity of simultaneity}
%%%%%%%%%%%%%%%%%%%%%%%%%%%%%%%%%%%%%%%%%%%%%%%%%%%%%%%%%%%%%%%%%%%%%%%%%%%

A rocket 4~ly long (in its own frame) moves at $0.6c$. A lamp at its exact
middle flashes at $t = t' = 0$, at the origin. Detectors at the back
($x' = -2$) and the front ($x'=+2$) record when the light reaches them.

\begin{center}
\begin{spacetime}[1.0]{-3}{6}{0}{6}
  \movinggrid{0.6}{0}{0}
  \answer{
    \draw[st ans] (-1.6,0) -- (2,6);
    \draw[st ans] (1.6,0) -- (5.2,6);
    \draw[st light] (0,0) -- (6,6);
    \draw[st light] (0,0) -- (-3,3);
    \node[st event, label={[st anslabel]left:$B$}] at (-1,1) {};
    \node[st event, label={[st anslabel]right:$F$}] at (4,4) {};
    \draw[st ans, dashed] (-1,1) -- (4,4);
  }
\end{spacetime}
\end{center}

\begin{questions}
  \question Trace the worldlines of the back and the front of the rocket
  along the tilted grid. Then draw the flash going each way.
  \begin{solution}[\room{1}]
    The ends follow the tilted lines $x'=-2$ and $x'=+2$. The flash is two
    $45^\circ$ lines from the origin.
  \end{solution}

  \question Mark $B$, where light reaches the back, and $F$, where it reaches
  the front. Give the rocket time $t'$ of each.
  \begin{solution}[\room{1}]
    $t'_B = t'_F = 2$~yr. In the rocket frame the two arrivals are simultaneous,
    as they must be: the lamp is in the middle and light goes the same speed both
    ways.
  \end{solution}

  \question Give the Earth time $t$ of each. Which arrival happens first in the
  Earth frame?
  \begin{solution}[\room{1}]
    $B$ at $(x,t)=(-1,1)$ and $F$ at $(x,t)=(4,4)$. In the Earth frame the light
    reaches the back first, 3~yr before it reaches the front.
  \end{solution}

\newpage
  \question Explain, from the point of view of someone on Earth, why the back
  is reached first.
  \begin{solution}[\room{1}]
    Light goes the same speed both ways in the Earth frame too, but the back of
    the rocket is moving toward the flash and the front is running away from it.
    The light going backward has less distance to cover.
  \end{solution}

  \question Connect $B$ and $F$. Which grid lines is your line parallel to?
  In one sentence, what does that say about the word ``simultaneous''?
  \begin{solution}[\room{1}]
    The line through $B$ and $F$ is parallel to the $x'$ axis: it is one of the
    rocket's ``now'' lines. Events that are simultaneous in one frame are not
    simultaneous in another moving relative to it; ``at the same time'' depends
    on the frame.
  \end{solution}
\end{questions}

\newpage
%%%%%%%%%%%%%%%%%%%%%%%%%%%%%%%%%%%%%%%%%%%%%%%%%%%%%%%%%%%%%%%%%%%%%%%%%%%
\section{Time dilation}
%%%%%%%%%%%%%%%%%%%%%%%%%%%%%%%%%%%%%%%%%%%%%%%%%%%%%%%%%%%%%%%%%%%%%%%%%%%

\begin{questions}
  \question A light clock is a pulse of light bouncing between two mirrors a
  distance $D$ apart; one tick is one round trip. In the clock's own frame a
  tick takes $\Delta t_0 = 2D/c$. Seen from a frame where the clock moves
  sideways at speed $v$, the light travels a longer, diagonal path. Use the
  Pythagorean theorem to show that the moving clock's tick takes
  \[ \Delta t = \gamma\,\Delta t_0, \qquad \gamma = \frac{1}{\sqrt{1-v^2/c^2}}. \]
  \begin{solution}[\room{1}]
    In the frame where the clock moves, each half-tick takes $\Delta t/2$. In
    that time the light goes $c\,\Delta t/2$ along the diagonal while the clock
    moves $v\,\Delta t/2$ sideways, and the vertical leg is still $D$:
    \[ \left(\tfrac{c\Delta t}{2}\right)^2 = D^2 + \left(\tfrac{v\Delta t}{2}\right)^2
       \quad\Rightarrow\quad
       \Delta t = \frac{2D/c}{\sqrt{1-v^2/c^2}} = \gamma\,\Delta t_0 . \]
  \end{solution}
\end{questions}

\newpage
The rocket carries a clock at $x' = 0$, and Earth has a clock at $x = 0$.
Both read zero at the origin.

\begin{center}
\begin{spacetime}[0.95]{-1}{6}{0}{8}
  \movinggrid{0.6}{0}{0}
  \answer{
    \draw[st ans] (0,0) -- (4.8,8);
    \node[st event, label={[st anslabel]right:$R$}] at (3,5) {};
    \node[st event, label={[st anslabel]left:$G$}] at (0,4) {};
    \draw[st ans, dashed] (0,4) -- (3.75,6.25);
    \draw[st ans, dashed] (3,5) -- (0,5);
  }
\end{spacetime}
\end{center}

\begin{questions}
  \question Mark the event $R$ where the rocket clock reads $t' = 4$~yr. How
  much time has passed in the Earth frame? What is the ratio, and how does it
  compare to $\gamma$?
  \begin{solution}[\room{1}]
    $R$ is at $(x,t)=(3,5)$. Earth says 5~yr have passed while the rocket clock
    shows 4: the ratio is $5/4 = 1.25 = \gamma$. Earth says the rocket clock
    runs slow.
  \end{solution}

  \question Mark the event $G$ where the Earth clock reads $t = 4$~yr. Count
  tilted lines: what is the rocket time $t'$ of $G$? What does the rocket
  conclude about the Earth clock?
  \begin{solution}[\room{1}]
    $G$ is at $(x',t') = (-3,5)$. The rocket says 5~yr have passed while the
    Earth clock shows 4. The rocket says the \emph{Earth} clock runs slow, by the
    same factor $\gamma$.
  \end{solution}

\newpage
  \question Each frame says the other's clock runs slow. Draw the rocket's
  ``now'' line through $G$, and Earth's ``now'' line through $R$. Use them to
  explain why the two claims do not contradict each other.
  \begin{solution}[\room{1}]
    Comparing two clocks that are far apart means deciding which reading on one
    happens ``at the same time'' as a reading on the other, and the frames
    disagree about that. Earth's now line through $R$ is horizontal, so it pairs
    rocket~4 with Earth~5. The rocket's now line through $G$ is parallel to $x'$
    and meets the rocket worldline at $t'=5$, so it pairs Earth~4 with rocket~5.
    Each comparison is right in its own frame; they are different comparisons.
  \end{solution}
\end{questions}

\newpage
%%%%%%%%%%%%%%%%%%%%%%%%%%%%%%%%%%%%%%%%%%%%%%%%%%%%%%%%%%%%%%%%%%%%%%%%%%%
\section{Length contraction: the pole and the barn}
%%%%%%%%%%%%%%%%%%%%%%%%%%%%%%%%%%%%%%%%%%%%%%%%%%%%%%%%%%%%%%%%%%%%%%%%%%%

A barn sits at rest on Earth, with its open doors at $x=0$ and $x=4$~ly. A
pole 5~ly long in its own frame moves through it at $0.6c$. The back of the
pole is at $x' = 0$ and the front at $x' = 5$, and the back reaches the first
door at the origin.

\begin{center}
\begin{spacetime}[0.95]{-1}{8}{-4}{4}
  \movinggrid{0.6}{0}{0}
  \answer{
    \draw[st ans] (0,-4) -- (0,4);
    \draw[st ans] (4,-4) -- (4,4);
    \draw[st ans] (-1,-1.6667) -- (3.4,5.6667);
    \draw[st ans] (1.6,-4) -- (8,6.6667);
    \node[st event, label={[st anslabel]above left:$P$}] at (0,0) {};
    \node[st event, label={[st anslabel]right:$Q$}] at (4,0) {};
    \draw[st ans, dashed] (0,0) -- (4,0);
    \draw[st ans, dashed] (4,0) -- (-1,-3); % the pole's now line t'=-3
  }
\end{spacetime}
\end{center}

\begin{questions}
  \question Draw the worldlines of both barn doors and of both ends of the pole.
  \begin{solution}[\room{1}]
    The doors are vertical lines at $x=0$ and $x=4$. The pole's ends follow the
    tilted lines $x'=0$ (the $t'$ axis) and $x'=5$.
  \end{solution}

  \question Earth measures the pole's length by marking both ends at the same
  Earth time, $t=0$. How long is the pole in the Earth frame? Compare to
  $5/\gamma$.
  \begin{solution}[\room{1}]
    At $t=0$ the back is at $x=0$ and the front at $x=4$, so Earth measures 4~ly,
    which is $5/1.25$. The moving pole is contracted.
  \end{solution}

  \question According to Earth, does the pole ever fit inside the barn?
  \begin{solution}[\room{1}]
    Yes, exactly, at $t=0$: back end at the first door (event $P$), front end at
    the far door (event $Q$).
  \end{solution}

\newpage
  \question Now take the pole's point of view. Count tilted lines to find the
  rocket time $t'$ of $P$ and of $Q$. Which happens first for the pole? In the
  pole's frame, how long is the barn?
  \begin{solution}[\room{1}]
    $t'_P = 0$ and $t'_Q = -3$~yr. For the pole, the front reaches the far door
    3~yr \emph{before} the back reaches the first door, so the pole never fits.
    The barn moves past the pole and is contracted to $4/1.25 = 3.2$~ly.
  \end{solution}

  \question Earth says the pole fits; the pole says it does not. Who is right?
  Use the diagram to explain.
  \begin{solution}[\room{1}]
    Both. ``The pole is inside the barn'' means both ends are inside \emph{at the
    same time}, and the two frames disagree about which events are simultaneous.
    Earth's now line $t=0$ pairs $P$ with $Q$; the pole's now lines are tilted
    and never pair them.
  \end{solution}
\end{questions}

\newpage
%%%%%%%%%%%%%%%%%%%%%%%%%%%%%%%%%%%%%%%%%%%%%%%%%%%%%%%%%%%%%%%%%%%%%%%%%%%
\section{The invariant interval}
%%%%%%%%%%%%%%%%%%%%%%%%%%%%%%%%%%%%%%%%%%%%%%%%%%%%%%%%%%%%%%%%%%%%%%%%%%%

The frames disagree about times, lengths and simultaneity. One combination of
coordinates they all agree on is the \emph{interval}, $s^2 = t^2 - x^2$.

\begin{questions}
  \question Fill in the table with events from earlier in the lab.

  \begin{center}
  \renewcommand{\arraystretch}{1.6}
  \begin{tabular}{l|c|c|c||c|c|c}
    Event & $x$ & $t$ & $t^2-x^2$ & $x'$ & $t'$ & $t'^2-x'^2$ \\ \hline
    $E$ (Section 1) & 3 & 5 &
      \ifprintanswers 16\fi & \ifprintanswers 0\fi & \ifprintanswers 4\fi & \ifprintanswers 16\fi \\
    $F$ (Section 1) & 5 & 5 &
      \ifprintanswers 0\fi & \ifprintanswers 2.5\fi & \ifprintanswers 2.5\fi & \ifprintanswers 0\fi \\
    $B$ (Section 2) & \ifprintanswers $-1$\fi & \ifprintanswers 1\fi &
      \ifprintanswers 0\fi & \ifprintanswers $-2$\fi & \ifprintanswers 2\fi & \ifprintanswers 0\fi \\
    $G$ (Section 3) & 0 & 4 &
      \ifprintanswers 16\fi & \ifprintanswers $-3$\fi & \ifprintanswers 5\fi & \ifprintanswers 16\fi \\
    $Q$ (Section 4) & 4 & 0 &
      \ifprintanswers $-16$\fi & \ifprintanswers 5\fi & \ifprintanswers $-3$\fi & \ifprintanswers $-16$\fi \\
  \end{tabular}
  \end{center}

  \question What do you notice? What does $s^2 = 0$ mean about an event and the
  origin?
  \begin{solution}[\room{1}]
    The interval comes out the same in both frames for every event. $s^2=0$ means
    the event lies on a light ray from the origin, and both frames agree on that.
  \end{solution}

\newpage
  \question On the grid below, plot every event you can find with
  $t^2 - x^2 = 16$ (try $x = 0, \pm 3, \pm 4, \pm 5$) and sketch the curve
  through them. Which tick on the rocket's $t'$ axis does the curve pass
  through?

  \begin{center}
  \begin{spacetime}[0.8]{-6}{6}{0}{8}
    \movinggrid{0.6}{0}{0}
    \answer{
      \draw[st ans, domain=-6:6, samples=60, smooth] plot (\x, {sqrt(16+\x*\x)});
      \foreach \p in {(0,4),(3,5),(-3,5),(4,5.657),(-4,5.657),(5,6.403),(-5,6.403)}
        {\node[st event] at \p {};}
    }
  \end{spacetime}
  \end{center}

  \begin{solution}[\room{1}]
    The curve is a hyperbola through $(0,4)$, $(\pm3, 5)$, $(\pm4, 5.66)$ and
    $(\pm5, 6.40)$. It passes through the $t' = 4$ tick.
  \end{solution}

  \question The interval from the origin to an event on a clock's worldline is
  the time that clock itself measures (its \emph{proper time}). Use that to
  explain why the tilted tick marks look farther apart than the square ones,
  even though each stands for 1~yr.
  \begin{solution}[\room{2}]
    Every clock that leaves the origin reads 4~yr when it reaches the
    $s^2=16$ hyperbola. The hyperbola curves upward and away from the origin, so
    along the tilted $t'$ axis it is farther out on the page than along the
    vertical $t$ axis. The page is drawn with Earth's ruler; distances on it do
    not measure the rocket's time.
  \end{solution}
\end{questions}

\newpage
%%%%%%%%%%%%%%%%%%%%%%%%%%%%%%%%%%%%%%%%%%%%%%%%%%%%%%%%%%%%%%%%%%%%%%%%%%%
\section{A real check: muons}
%%%%%%%%%%%%%%%%%%%%%%%%%%%%%%%%%%%%%%%%%%%%%%%%%%%%%%%%%%%%%%%%%%%%%%%%%%%

Muons are made when cosmic rays hit the top of the atmosphere. At rest, a muon
lives an average of $\tau = 2.2~\mu$s. In 1962, David Frisch and James Smith
counted muons moving at about $0.995c$ at the top of Mount Washington, 1917~m
up, and then at sea level in Cambridge, Massachusetts.\footnote{D.~H.~Frisch and
J.~H.~Smith, ``Measurement of the relativistic time dilation using
$\mu$-mesons,'' \emph{Am.\ J.\ Phys.}\ \textbf{31}, 342 (1963).}

\begin{center}
\begin{tabular}{lc}
  \toprule
  Location & Muons per hour \\ \midrule
  Mount Washington, 1917~m & 563 \\
  Cambridge, sea level     & 412 \\ \bottomrule
\end{tabular}
\end{center}

For muons it helps to keep $c=1$ but change units: distance in
light-microseconds (1~light-$\mu$s $= 300$~m) and time in $\mu$s.

\begin{questions}
  \question How far is the trip down, in light-$\mu$s? How long does it take,
  in the Earth frame?
  \begin{solution}[\room{1}]
    $1917~\text{m} / 300~\text{m} = 6.39$~light-$\mu$s. At $0.995c$ the trip
    takes $6.39/0.995 = 6.4~\mu$s.
  \end{solution}

  \question If a muon's clock ran at the Earth rate, the fraction surviving the
  trip would be $e^{-\Delta t/\tau}$. How many of the 563 per hour should reach
  sea level? Compare to what was measured.
  \begin{solution}[\room{1}]
    $e^{-6.4/2.2} = e^{-2.9} \approx 0.054$, so about 30 per hour (Frisch and
    Smith, correcting for the spread in speeds, predicted 27). They measured 412,
    more than ten times as many.
  \end{solution}

  \question From the measured counts, how much time passed on the muons' own
  clocks during the trip? What time-dilation factor does that imply? Compare to
  $\gamma$ at $0.995c$.
  \begin{solution}[\room{1.3}]
    $\Delta t_0 = \tau\ln(563/412) = 2.2 \times 0.312 = 0.69~\mu$s. The implied
    factor is $6.4/0.69 \approx 9$. At $0.995c$, $\gamma = 10.0$. Close: the
    muons slow down in the air on the way, so their average $\gamma$ is a little
    lower. (Frisch and Smith's careful analysis gave $8.8 \pm 0.8$.)
  \end{solution}

\newpage
  \question In the muon's own frame the muon is at rest and lives its normal
  life. How does the muon explain making it down? How tall is the mountain in
  the muon frame?
  \begin{solution}[\room{1}]
    In the muon frame the mountain rushes up at $0.995c$ and is length-contracted
    by $\gamma$: about $1917/9 \approx 210$~m, which takes only $0.7~\mu$s to pass.
    Earth says time dilation, the muon says length contraction: one fact, two
    descriptions.
  \end{solution}

  \question Sketch a spacetime diagram of the trip in the Earth frame: the
  mountaintop and sea level as worldlines, and the muon's worldline between
  them. Label the Earth time and the muon's own time for the trip. (The grids
  in this lab are for $0.6c$, so this one is a sketch.)
  \begin{solution}[\room{2}]
    Two vertical lines, the top at $x=0$ and sea level at $x = 6.4$~light-$\mu$s
    (taking downward as $+x$). The muon's worldline runs from the top at $t=0$
    to sea level at $t=6.4~\mu$s, very nearly at $45^\circ$. Its own clock reads
    only about $0.7~\mu$s at the bottom.
  \end{solution}
\end{questions}

\newpage
%%%%%%%%%%%%%%%%%%%%%%%%%%%%%%%%%%%%%%%%%%%%%%%%%%%%%%%%%%%%%%%%%%%%%%%%%%%
\section{The twin paradox}
%%%%%%%%%%%%%%%%%%%%%%%%%%%%%%%%%%%%%%%%%%%%%%%%%%%%%%%%%%%%%%%%%%%%%%%%%%%

One twin stays on Earth at $x=0$. The other leaves at $t=0$, flies at $0.6c$ to
a star 3~ly away, turns around immediately, and flies home at $0.6c$. Both
diagrams show the whole trip. The left one carries the outbound rocket's grid
$(x',t')$, with its origin where the trip starts. The right one carries the
homebound rocket's grid $(x'',t'')$, tilted the other way, with its origin at
the turnaround.

\begin{center}
\begin{spacetime}[0.9]{-1}{5}{0}{10}
  \movinggrid{0.6}{0}{0}
  \answer{
    \draw[st ans] (0,0) -- (0,10);
    \draw[st ans] (0,0) -- (3,5) -- (0,10);
    \node[st event, label={[st anslabel]right:$T$}] at (3,5) {};
    \draw[st ans, dashed] (3,5) -- (0,3.2);
  }
\end{spacetime}
\hspace{1.2cm}
\begin{spacetime}[0.9]{-1}{5}{0}{10}
  \movinggrid['']{-0.6}{3}{5}
  \answer{
    \draw[st ans] (0,0) -- (0,10);
    \draw[st ans] (0,0) -- (3,5) -- (0,10);
    \node[st event, label={[st anslabel]right:$T$}] at (3,5) {};
    \draw[st ans, dashed] (3,5) -- (0,6.8);
  }
\end{spacetime}
\end{center}

\begin{questions}
  \question On both diagrams, draw both twins' worldlines and mark the
  turnaround $T$. How much does the Earth twin age?
  \begin{solution}[\room{1}]
    The trip covers 6~ly at $0.6c$, so the Earth twin ages 10~yr. The traveler's
    worldline is the tilted line out to $T=(3,5)$ and back to $(0,10)$.
  \end{solution}

  \question Use the left grid to read the traveler's clock at $T$, and the
  right grid to read how much it advances on the way home. How much does the
  traveler age? Check your answer with the interval.
  \begin{solution}[\room{1}]
    4~yr out ($t'=4$ at $T$) and 4~yr back ($t''$ goes from 0 to 4). The traveler
    ages 8~yr, two fewer than the Earth twin. Check: each leg has
    $\sqrt{5^2-3^2} = 4$.
  \end{solution}

\newpage
  \question On each leg, the traveler says the Earth twin's clock runs slow.
  So why isn't the traveler the older one? What makes the two twins' situations
  different? Point to it on the diagram.
  \begin{solution}[\room{1}]
    The traveler's worldline has a kink at $T$: the traveler changes frames,
    and the Earth twin never does. The slow-clock argument holds within one
    frame, and the traveler uses two of them.
  \end{solution}

  \question Draw the traveler's ``now'' line through $T$ on each diagram, and
  read where it meets the Earth twin's worldline. What happens to the
  traveler's idea of ``what time it is on Earth'' at the turnaround?
  \begin{solution}[\room{1.3}]
    Just before turning, the outbound now line through $T$ meets Earth at
    $t = 3.2$~yr. Just after, the homebound now line meets Earth at $t = 6.8$~yr.
    At the turnaround the traveler's ``now on Earth'' jumps ahead 3.6~yr. Earth
    ages $3.2$ (slow, outbound) $+\ 3.6$ (the jump) $+\ 3.2$ (slow, homebound)
    $= 10$~yr, as it should.
  \end{solution}
\end{questions}

\newpage
%%%%%%%%%%%%%%%%%%%%%%%%%%%%%%%%%%%%%%%%%%%%%%%%%%%%%%%%%%%%%%%%%%%%%%%%%%%
% \section{Extension (optional): adding velocities}
% %%%%%%%%%%%%%%%%%%%%%%%%%%%%%%%%%%%%%%%%%%%%%%%%%%%%%%%%%%%%%%%%%%%%%%%%%%%

% The rocket, moving at $0.6c$, fires a probe forward at the origin. In the rocket
% frame the probe moves at $0.6c$.

% \begin{center}
% \begin{spacetime}[0.8]{0}{9}{0}{10}
%   \movinggrid{0.6}{0}{0}
%   \answer{
%     \draw[st ans] (0,0) -- (8.8235,10);
%     \node[st event, label={[st anslabel]left:$P$}] at (7.5,8.5) {};
%   }
% \end{spacetime}
% \end{center}

% \begin{questions}
%   \question In the rocket frame, where is the probe at $t' = 5$~yr? Mark that
%   event $P$ using the tilted grid, and draw the probe's worldline.
%   \begin{solution}[\room{0.6}]
%     At $(x',t') = (3,5)$, which is $(x,t)=(7.5, 8.5)$ in the Earth frame.
%   \end{solution}

%   \question Read the Earth coordinates of $P$ and find the probe's speed in the
%   Earth frame. What would you have predicted before this lab?
%   \begin{solution}[\room{1}]
%     $7.5/8.5 = 0.88c$. Simple addition says $0.6c + 0.6c = 1.2c$, faster than
%     light. That cannot be right, because both frames agree on the speed of light.
%   \end{solution}

%   \question The relativistic rule is
%   $\displaystyle u = \frac{u' + v}{1 + u'v/c^2}$. Check it against your
%   reading.
%   \begin{solution}[\room{1}]
%     $(0.6+0.6)/(1+0.36) = 1.2/1.36 = 0.882$, matching the diagram.
%   \end{solution}
% \end{questions}

\end{document}
